% Adapted from nested-list-mutation-with-diagram.tex for sequence operations before mutation.
\question What would Python display? Run the statements in order.
For nested expressions, work from the innermost expression outward.

\begin{blocksection}
\begin{lstlisting}
>>> a = [1, 2, 3]
>>> a
\end{lstlisting}
\begin{solution}[.25in]
\begin{lstlisting}
[1, 2, 3]
\end{lstlisting}
\end{solution}
\end{blocksection}

\begin{blocksection}
\begin{lstlisting}
>>> a[2]
\end{lstlisting}
\begin{solution}[.25in]
\begin{lstlisting}
3
\end{lstlisting}
\end{solution}
\end{blocksection}

\begin{blocksection}
\begin{lstlisting}
>>> a[-1]
\end{lstlisting}
\begin{solution}[.25in]
\begin{lstlisting}
3
\end{lstlisting}
\end{solution}
\end{blocksection}

\begin{blocksection}
\begin{lstlisting}
>>> a[1:]
\end{lstlisting}
\begin{solution}[.25in]
\begin{lstlisting}
[2, 3]
\end{lstlisting}
\end{solution}
\end{blocksection}

\begin{blocksection}
\begin{lstlisting}
>>> a[:2]
\end{lstlisting}
\begin{solution}[.25in]
\begin{lstlisting}
[1, 2]
\end{lstlisting}
\end{solution}
\end{blocksection}

\begin{blocksection}
\begin{lstlisting}
>>> a + [4, 5]
\end{lstlisting}
\begin{solution}[.25in]
\begin{lstlisting}
[1, 2, 3, 4, 5]
\end{lstlisting}
\end{solution}
\end{blocksection}

\begin{blocksection}
\begin{lstlisting}
>>> a
\end{lstlisting}
\begin{solution}[.25in]
\begin{lstlisting}
[1, 2, 3]
\end{lstlisting}
\end{solution}
\end{blocksection}

\begin{blocksection}
\begin{lstlisting}
>>> nested = [1, 2, 3, 4, [5, 6]]
>>> nested[4]
\end{lstlisting}
\begin{solution}[.25in]
\begin{lstlisting}
[5, 6]
\end{lstlisting}
\end{solution}
\end{blocksection}

\begin{blocksection}
\begin{lstlisting}
>>> nested[4][0]
\end{lstlisting}
\begin{solution}[.25in]
\begin{lstlisting}
5
\end{lstlisting}
\end{solution}
\end{blocksection}

\begin{blocksection}
\begin{lstlisting}
>>> nested[3:5]
\end{lstlisting}
\begin{solution}[.25in]
\begin{lstlisting}
[4, [5, 6]]
\end{lstlisting}
\end{solution}
\end{blocksection}

\begin{questionmeta}
\textbf{Teaching Tips}\par\nopagebreak[4]
Suggested Time: 6--8 min; Difficulty: Easy
\nopagebreak[4]
\begin{itemize}
    \item \textbf{Indexing:} Indices start at zero; negative indices count from the end. Ask students to predict the type of each result before its value.
    \item \textbf{Slicing:} Include the start index and exclude the stop index. An omitted bound extends to the corresponding end of the list.
    \item \textbf{Concatenation:} The expression \lstinline{a + [4, 5]} produces a list value; evaluating it does not change the value of \lstinline{a}.
    \item \textbf{Nested access:} First evaluate \lstinline{nested[4]}, then select index zero from that result. Focus on the elements and the operation at each step.
\end{itemize}
\end{questionmeta}
